\documentclass[10pt,a4paper]{amsart}
\usepackage[margin=19mm]{geometry}
\usepackage{amsmath,amssymb,mathtools}
\usepackage[scr=rsfs,cal=euler]{mathalfa}
\usepackage{xfrac}
\usepackage[unicode,hidelinks,bookmarks=false]{hyperref}
\newcommand{\ie}{i.e.\ }
\newcommand{\cf}{cf.\ }
\newcommand{\resp}{respectively\ }
\newcommand{\Subsection}{\subsection}
\providecommand{\nobreakdash}{\nobreak\hbox{-}}
\setlength{\emergencystretch}{2em}
\multlinegap0pt
\usepackage{bm}
\usepackage{bbm}
\usepackage{dsfont}
\usepackage{xspace}
\usepackage{cancel}
\usepackage{mathtools}
\usepackage{amsthm}
\makeatletter
\@ifundefined{theorem}{\newtheorem{theorem}{Theorem}}{}
\@ifundefined{lemma}{\newtheorem{lemma}[theorem]{Lemma}}{}
\makeatother
\title[Stranding $\mathfrak{sl}\_n$ webs]{Stranding $\mathfrak{sl}\_n$ webs}
\author{Heather M. Russell, Julianna Tymoczko}
\date{Source: arXiv:2510.12035v3 (CC BY 4.0)}
\begin{document}
\maketitle

\noindent\emph{Source and licence.} Stranding $\mathfrak{sl}_n$ webs. Version arXiv:2510.12035v3 (CC BY 4.0), distributed under the Creative Commons Attribution 4.0 International licence, as recorded in the arXiv licence metadata for this exact version and shown on the abstract page https://arxiv.org/abs/2510.12035v3. Full rights notice: licenses/arxiv-2510.12035-CC-BY-4.0.txt.\par\smallskip

\noindent\textit{Editorial scope.} Counted unit: the Theorem whose source label is \texttt{thm: strandings are bijective with binary labelings} in the article (source0.tex lines 1338--1349, source label \texttt{thm: strandings are bijective with binary labelings}), delivered together with the article's own complete original proof of it (source0.tex lines 1351--1384, one \texttt{proof} environment of 34 lines). No mathematics is written, repaired or re-derived: statement and proof are the article's own text line for line. Required context retained uncounted: thm:net zero strand contribution at each vertex (lines 959--964); lem:second validity condition for strandings in terms of binary labels (lines 1185--1194); equation: defining b\_S(e) (lines 1229--1232). Every definition, display and label of those lines is retained unchanged. Omitted from this package and not redistributed: 1--958, 965--1184, 1195--1228, 1233--1337, 1350--1350, 1385--5966. Nothing inside a retained excerpt is omitted or altered: every retained body is delivered exactly as the article's lines are written, so no omission receipt of any kind accompanies this package. Citations: the retained excerpts carry no citation command at all, so no bibliography entry of the article is transcribed and no reference is redistributed. Reference adaptation: none was needed and none was made; every reference inside the delivered unit resolves inside it, so no unresolved reference or citation is produced. Printed slips kept literal: no printed word, symbol, hypothesis or number of the counted statement or of its proof was altered. Layout and class: the article's own class is \texttt{amsart} with packages including geometry, latexsym, amsfonts, amssymb, amstext, mathrsfs, framed, youngtab, graphicx, tikz, subcaption, tikz-cd, hyperref, inputenc; the wrapper is the lane's pinned \texttt{amsart} editorial wrapper at 10pt on A4, which restates the theorem-like environments the retained text needs and adds the article's own macro definitions that the retained text needs (none), with extra wrapper packages (bm, bbm, dsfont, xspace, cancel, mathtools). The delivered numbering is the unit's own. Assets: no retained span contains a figure environment, a graphics-inclusion command, a PDF-page inclusion, a table or a TikZ picture; no image file is copied into this package, the package declares no asset dependency and no image file is redistributed with it. Licence: version arXiv:2510.12035v3 of this article is distributed under the Creative Commons Attribution 4.0 International licence, as recorded in the arXiv licence metadata for this exact version and shown on the abstract page https://arxiv.org/abs/2510.12035v3. A full rights notice is delivered at licenses/arxiv-2510.12035-CC-BY-4.0.txt. Characters: every retained excerpt is pure ASCII, so the delivered engine, XeLaTeX with its default Latin Modern text font, reports no missing character.
\begin{theorem}\label{thm:net zero strand contribution at each vertex}
Let $G$ be a web graph. For each edge $e\in E(G)$, choose valid edgewise simple strand data $S(e)=S_+(e)\sqcup S_-(e)$. These choices yield a valid stranding of $G$ if and only if for each interior vertex $v$ with incident edges $e_1,e_2,e_3$ and each color $1\le c\le n-1$,
$$
\sum_{i=1}^3 \sigma_v(e_i)\alpha_c^\vee(e_i)=0.
$$
\end{theorem}

\begin{lemma}\label{lem:second validity condition for strandings in terms of binary labels}
Let $\vec{b}$ be a binary vector with exactly $\ell$ ones, where $1\le \ell\le n-1$. Then
$$
\sum_{i=1}^{n-1}\alpha_i^\vee(\vec{b})\,i=
\begin{cases}
\ell & \textup{if } b_n=0,\\
\ell-n & \textup{if } b_n=1.
\end{cases}
$$
\end{lemma}

\begin{equation}\label{equation: defining b_S(e)}
\alpha_c^\vee(b_S(e))=b_S(e)_c-b_S(e)_{c+1}=\alpha_c^\vee(e)
\qquad\textup{for } 1\le c\le n-1.
\end{equation}

\begin{theorem}\label{thm: strandings are bijective with binary labelings}
For any $\mathfrak{sl}_n$ web graph $G$, the assignments
$$
b\longmapsto S_b,
\qquad
S\longmapsto b_S
$$
define a bijection
$$
\mathcal{B}in(G)\cong \mathcal{S}tr(G).
$$
\end{theorem}

\begin{proof}
We first check the constructions edgewise. Let $e$ be an edge of $G$.

Let $b\in\mathcal{B}in(G)$. By definition, the simple strands of $S_b(e)$ are determined by the signs of the differences $\alpha_c^\vee(b(e))=b_c-b_{c+1}$. These nonzero values necessarily alternate in sign, so the first validity condition holds. The second validity condition follows from Lemma~\ref{lem:second validity condition for strandings in terms of binary labels}.

Conversely, let $S\in\mathcal{S}tr(G)$. Equation~\eqref{equation: defining b_S(e)} determines $b_S(e)$ uniquely from the values $\alpha_c^\vee(e)$ together with the value of the last entry. Because the nonzero values $\alpha_c^\vee(e)$ alternate in sign, the vector $b_S(e)$ is binary. The second validity condition for $S$ implies that $b_S(e)$ has exactly as many ones as the weight of $e$.


It remains to check the vertex conditions. Let $v$ be an interior vertex with incident edges $e_1,e_2,e_3$. 

Suppose $S\in\mathcal{S}tr(G)$. By Theorem~\ref{thm:net zero strand contribution at each vertex},
$$
\sum_{i=1}^3 \sigma_v(e_i)\alpha_c^\vee(e_i)=0
\qquad\textup{for all } c.
$$
Applying the coroots $\alpha_c^\vee$ to the vector
$$
\sum_{i=1}^3 \sigma_v(e_i)b_S(e_i)
$$
therefore gives zero for every $c$, so this vector is a multiple of $\vec{1}$. Hence $b_S$ is a binary labeling.

Now let $b\in\mathcal{B}in(G)$.  By the definition of binary labeling, we have
$$
\sum_{i=1}^3 \sigma_v(e_i)b(e_i)\in \mathbb{Z}\vec{1}.
$$
 Therefore, applying $\alpha_c^\vee$ to this expression yields
$$
\sum_{i=1}^3 \sigma_v(e_i)\alpha_c^\vee(b(e_i))=0
\qquad\textup{for all } c.
$$
Since $\alpha_c^\vee(b(e_i))=\alpha_c^\vee(e_i)$ by construction of the stranding $S_b$, Theorem~\ref{thm:net zero strand contribution at each vertex} shows that the edgewise strand data glue to a valid stranding. Thus $S_b\in\mathcal{S}tr(G)$.

Finally, a straightforward calculation shows $(b_{S_b}(e))_n=b(e)_n$ for each edge $e$ of $G$. Since the two constructions recover the same values of all $\alpha_c^\vee$ on every edge, we conclude they are inverse to one another.
\end{proof}

\end{document}
